\documentclass[a4paper,10pt]{article}
\usepackage[margin=1.55cm]{geometry}
\usepackage[dutch]{babel}
\usepackage[T1]{fontenc}
\usepackage{lmodern}
\usepackage{enumitem,titlesec,xcolor}
\usepackage[hidelinks]{hyperref}
\definecolor{accent}{HTML}{173B50}
\pagestyle{empty}
\setlength{\parindent}{0pt}
\setlength{\parskip}{3pt}
\titleformat{\section}{\large\bfseries\color{accent}}{}{0pt}{}[\titlerule]
\titlespacing*{\section}{0pt}{9pt}{4pt}
\setlist[itemize]{leftmargin=14pt,nosep,topsep=2pt}
\newcommand{\entry}[3]{\textbf{#1}\hfill #2\\#3\par}
\begin{document}
{\LARGE\bfseries\color{accent} Bradley Velle}\\[3pt]
{\large Student Toegepaste Informatica | Softwareontwikkeling en procesanalyse}

Eernegem, België \enspace | \enspace \href{tel:+32468409996}{+32 468 40 99 96} \enspace | \enspace \href{mailto:info@velleb.com}{info@velleb.com}\\
\href{https://velleb.com}{velleb.com} \enspace | \enspace \href{https://github.com/Br4dleyV}{github.com/Br4dleyV} \enspace | \enspace \href{https://linkedin.com/in/bradley-velle}{linkedin.com/in/bradley-velle}

\section{Profiel}
Derdejaarsstudent Toegepaste Informatica met projectervaring in webontwikkeling en procesanalyse. Werkt met Domain Storytelling en Event Storming aan de voorbereiding van een herbouw in .NET. Bouwde een eigen portfolio met Next.js, React en TypeScript. Wil doorgroeien in applicatiearchitectuur en de samenhang tussen software, data en bedrijfsprocessen.

\section{Opleiding}
\entry{Bachelor Toegepaste Informatica, Howest, Brugge}{2024--2027}{Lopend, derde jaar. Verwachte afstudeerdatum: 2027.}
\entry{Opleiding Graduaat Programmeren, Howest, Brugge}{2023--2024}{Opleiding gevolgd, daarna overgestapt naar de bachelor Toegepaste Informatica.}
\entry{Economie--Moderne Talen, SIGO, Eernegem/Gistel}{2017--2023}{Secundair onderwijs.}

\section{Relevante projecten}
\textbf{Teamproject voor de Cel Vermiste Personen, Belgische politie} \hfill 2026--heden
\begin{itemize}
\item Met het team van 4 studenten processen en domeinkennis in kaart brengen via Domain Storytelling en Event Storming.
\item De herbouw voorbereiden van een verouderde toepassing op basis van SharePoint en Access naar .NET, met schaalbaarheid als aandachtspunt. Het project bevindt zich in de analysefase.
\end{itemize}
\textbf{Eigen portfolio en blog: \href{https://velleb.com}{velleb.com}} \hfill Next.js, React, TypeScript
\begin{itemize}
\item Bouwde een responsieve website met Tailwind CSS, een Nederlandstalige en Engelstalige weergave en een licht en donker thema.
\item Richtte een blog in met Markdown-bestanden, serververwerking en expliciete publicatie van concepten, voegde zoeken en een afdrukweergave voor pdf-export toe.
\end{itemize}

\section{Technische vaardigheden}
\textbf{Talen:} C\#, TypeScript, JavaScript, Python, Java, SQL, HTML/CSS.\\
\textbf{Web en frameworks:} .NET / ASP.NET Core, React, Next.js, Vue, Nuxt, Node.js, Express, Tailwind CSS.\\
\textbf{Databases en API's:} PostgreSQL, MySQL, MariaDB, SQL Server, MongoDB, Redis, GraphQL.\\
\textbf{Cloud en infrastructuur:} Docker, Kubernetes, Microsoft Azure, Terraform, Ubuntu.\\
\textbf{Ontwikkeltools en platformen:} Git, GitHub, GitLab, Vercel, Cloudflare, Supabase.\\
\textbf{Analyse:} Domain Storytelling, Event Storming (lopend teamproject).

\section{Werkervaring}
\entry{Studentenjobs, Proxy Delhaize, Jumbo en VIVES-boekenverkoop}{2022--heden}{Klanten geholpen, kassawerk en administratieve taken uitgevoerd.}
\entry{Student documentbeheer, iCount, Westkerke}{2023}{Boekhoudkundige documenten gedigitaliseerd en geordend voor digitale archivering, geholpen bij de voorbereiding van btw-dossiers.}
\entry{Student IT en technische dienst, Hogeschool VIVES, Brugge}{mrt.--aug. 2022}{Computers opnieuw geïnstalleerd en geconfigureerd na een beveiligingsincident, nieuwe IT-apparatuur uitgerold en geholpen bij de inrichting van een nieuw campusgebouw.}

\section{Talen en mobiliteit}
\textbf{Nederlands:} moedertaal \enspace | \enspace \textbf{Engels:} vloeiend \enspace | \enspace \textbf{Frans en Duits:} basiskennis.\\
\textbf{Rijbewijzen:} B en A2.
\end{document}
